\clearpage
\section[\MakeUppercase{Обзор методов решения задачи 3D Bin Packing}]{\centering\MakeUppercase{Обзор методов решения задачи 3D Bin Packing}}

\subsection{\centering{Классические эвристики}}

Для задач упаковки часто применяются эвристики последовательного размещения: First Fit, Best Fit, First Fit Decreasing и их двумерные или трехмерные аналоги~\cite{coffman1997,lodi2002}. Общая идея таких методов состоит в том, что объекты упорядочиваются по некоторому приоритету, после чего каждый объект размещается в первую или лучшую допустимую позицию.

В прикладной задаче паллетизации порядок размещения становится самостоятельным параметром качества. Сортировка только по объему обычно повышает плотность нижних слоев, но может ухудшать обращение с хрупкими товарами. Сортировка по массе помогает удерживать тяжелые товары внизу, однако иногда приводит к тому, что легкие, но крупные короба оказываются неразмещенными. Поэтому практические эвристики часто используют составной приоритет: объем, массу, высоту, хрупкость, разрешенные ориентации и наличие запрета на штабелирование.

Отличие логистической постановки от абстрактной задачи упаковки состоит в том, что допустимое решение должно быть удобно не только для математической модели, но и для последующей перевозки. В задачах загрузки контейнеров и совмещенного планирования маршрутов также подчеркивается, что геометрическая плотность должна рассматриваться совместно с ограничениями реального обращения с грузом~\cite{gendreau2006}. Для рассматриваемого проекта это означает, что метод не может оцениваться только по заполненному объему: важны полнота размещения, устойчивость, хрупкость и время расчета.

В трехмерной постановке важную роль играют кандидатные позиции. Полный перебор всех координат невозможен, поэтому обычно используются специальные множества точек: углы свободных областей, extreme points, границы уровней или прямоугольники свободного пространства. Такие методы дают быстрое допустимое решение и хорошо подходят как baseline.

Недостаток последовательных эвристик состоит в их локальности. Если ранний выбор оказался неудачным, алгоритм обычно не возвращается к нему и не перестраивает уже созданную укладку. В задачах паллетизации это особенно заметно: один крупный короб может заблокировать полезную область, а хрупкий или нештабелируемый товар на нижнем уровне может испортить всю структуру над собой.

Несмотря на это, классические эвристики остаются необходимой частью экспериментального контура. Во-первых, они дают понятный baseline, с которым можно сравнивать более сложные методы. Во-вторых, быстрые эвристики часто используются как начальная фаза для метаэвристик: хорошее стартовое решение уменьшает объем последующего поиска. В-третьих, простые методы позволяют быстро выявлять ошибки в валидаторе и формате данных, так как их поведение легко интерпретировать вручную.

\subsection{\centering{Extreme Points и Bottom-Left-Fill}}

Подход extreme points строит набор потенциальных точек размещения на основе уже поставленных коробов. После размещения очередного объекта появляются новые точки около его правой, дальней и верхней граней. Для каждой точки перебираются допустимые ориентации, затем проверяются границы паллеты, пересечения, опора и дополнительные ограничения. Эвристики на базе extreme points специально разработаны для трехмерной упаковки и позволяют резко сократить число рассматриваемых координат без перехода к грубой дискретизации пространства~\cite{crainic2008}.

В исследуемом проекте greedy-решатель использует extreme points вместе с gravity drop. Фактически кандидатная точка задает координаты $x$ и $y$, а координата $z$ выбирается как минимальная высота, на которую короб может быть опущен без пересечения и с учетом штабелируемости нижних объектов. Такое поведение близко к эвристике Bottom-Left-Fill: предпочтение отдается более низким, левым и ближним позициям.

Практический смысл gravity drop состоит в том, что кандидатная точка не рассматривается как окончательная высота размещения. Если короб можно опустить ниже без пересечения с уже размещенными объектами, алгоритм делает это и тем самым уменьшает число висящих или неестественно высоких позиций. Для паллетизации это важно: низкий центр массы и заполнение нижних уровней обычно повышают устойчивость всей конструкции.

Набор extreme points после каждого размещения должен очищаться от заведомо бесполезных точек. В противном случае число кандидатов быстро растет, а решатель тратит большую часть времени на повторную проверку координат, которые уже заняты или находятся вне допустимой области. Типичные фильтры удаляют дубликаты, точки внутри уже размещенных коробов, точки за пределами паллеты и позиции, для которых невозможно поставить ни одну допустимую ориентацию. Такая фильтрация не меняет постановку задачи, но существенно влияет на скорость базового greedy и на качество начальных решений LNS.

Ограничение подхода extreme points связано с тем, что он сохраняет только локально естественные позиции. Если оптимальное решение требует временно оставить свободное место или разместить короб не в ближайшей нижней позиции, жадный декодер может не найти такую конфигурацию. Поэтому extreme points хорошо подходят для быстрого построения решений, но в сложных случаях требуют надстройки в виде локального поиска, рестартов или крупноокрестностной перестройки.

\subsection{\centering{MaxRects и послойная укладка}}

Алгоритмы семейства MaxRects широко используются в двумерной упаковке прямоугольников. Они поддерживают список свободных прямоугольных областей и после размещения объекта разбивают занятое пространство на новые свободные области. В трехмерной паллетизации такой подход можно применять послойно: выбрать высоту слоя, упаковать в нем прямоугольные основания коробов, затем перейти к следующему уровню.

В проекте файл \code{solve\_conditions.py} реализует MaxRects-подобный послойный решатель, адаптированный из идеи соревнования Packing Santa's Sleigh~\cite{kaggle_sleigh}. Метод быстро строит плотные слои и показывает высокую утилизацию объема, однако из-за послойной природы и сортировки по весу хуже контролирует хрупкость: на контрольном наборе средний fragility score равен 0.5198.

Главное преимущество послойной схемы заключается в уменьшении размерности промежуточной задачи. Вместо полного трехмерного поиска решатель фиксирует высоту слоя и решает двумерную задачу размещения оснований. Это делает метод быстрым и хорошо объяснимым: в каждом слое можно анализировать свободные прямоугольники, правила выбора позиции и причины отказа от размещения конкретного товара. Для инженерного прототипа такая прозрачность важна, поскольку позволяет быстро сравнивать изменения в эвристике.

Слабое место MaxRects-подобной схемы -- выбор высоты слоя. Если слой выбран по одному крупному коробу, рядом с ним могут оказаться объекты меньшей высоты, и над ними образуется неиспользуемый вертикальный объем. Если же слои часто менять, повышается фрагментация и появляются узкие свободные области, которые трудно заполнить последующими коробами. В реальной паллетизации эта проблема усиливается ограничениями \code{strict\_upright} и \code{stackable}: не каждый товар можно повернуть или использовать как опору для следующего уровня.

Для рассматриваемого проекта MaxRects играет роль сильного быстрого метода. Он полезен не только как самостоятельный решатель, но и как ориентир для интерпретации результатов: если другой алгоритм проигрывает MaxRects по плотности, но выигрывает по хрупкости и полноте заказа, это показывает реальный компромисс между геометрическим заполнением и логистической пригодностью укладки.

\subsection{\centering{Large Neighborhood Search}}

Large Neighborhood Search относится к метаэвристикам крупноокрестностного поиска~\cite{pisinger2010}. В отличие от малых локальных перестановок, LNS на каждой итерации разрушает существенную часть текущего решения, а затем восстанавливает ее. Такой подход особенно полезен для задач упаковки, где качество сильно зависит от взаимного расположения групп объектов.

Типовой цикл LNS состоит из следующих этапов:
\begin{enumerate}
  \item Построение начального решения;
  \item Выбор фрагмента решения для удаления;
  \item Восстановление частичного решения;
  \item Оценка нового решения;
  \item Принятие или отклонение нового состояния.
\end{enumerate}

В исследуемой реализации LNS начинается с усиленной greedy-фазы, затем применяет несколько destroy-эвристик и beam search на этапе repair. Критерий принятия основан на simulated annealing, что позволяет иногда принимать ухудшения и выходить из локальных максимумов.

Смысл destroy-фазы состоит в том, чтобы временно отказаться от части уже построенной укладки. Удаление может быть случайным, направленным на верхние уровни, связанным с локальным пространственным кластером или ориентированным на проблемные элементы, например на короба около нарушений хрупкости. Разные варианты разрушения дают разные типы окрестностей: одни позволяют улучшить плотность, другие исправляют плохую структуру поддержки, третьи помогают найти место для ранее неразмещенных товаров.

Repair-фаза восстанавливает решение, но уже в измененном контексте. Если удаленные короба вставляются в прежнем порядке и по тем же правилам, метод почти не отличается от повторного greedy. Поэтому в проекте используется beam search: на каждом шаге сохраняется несколько перспективных частичных вариантов, а не только один лучший. Это увеличивает вычислительные затраты, но позволяет не отбрасывать решение из-за локально спорного выбора, который может оказаться выгодным через несколько последующих размещений.

Принятие кандидата по simulated annealing добавляет управляемую стохастичность. Если принимать только улучшения, поиск быстро стабилизируется и перестает выходить из локального максимума. Если принимать ухудшения слишком часто, алгоритм теряет направленность и превращается в случайный перебор. Температура и скорость охлаждения задают баланс между этими режимами: в начале поиска допустимы более смелые перестройки, ближе к концу предпочтение отдается улучшению найденного лучшего решения.

Для задачи паллетизации LNS особенно удобен тем, что сохраняет валидатор как внешний модуль. Алгоритм может экспериментировать с разрушением и восстановлением, но итоговая оценка всегда проходит через единые правила: границы, пересечения, массу, опору, ориентации и хрупкость. Это снижает риск, что более сложный метод выиграет только за счет неявного ослабления ограничений.

\subsection{\centering{Генетические алгоритмы и ML-подходы}}

Генетические алгоритмы применяются к 3D Bin Packing через кодирование порядка размещения, ориентаций или параметров декодера. Особь обычно не хранит координаты напрямую, а задает правила, по которым greedy-декодер строит укладку. Это уменьшает риск невалидных решений и позволяет использовать эволюционный поиск в пространстве перестановок.

Современные исследования также рассматривают гибридные схемы, где эволюционный поиск дополняется генеративными моделями~\cite{zhang2024}. В таком подходе генератор пытается порождать перспективные перестановки или признаки решений, а дискриминатор помогает отличать их от слабых конфигураций. В проекте реализован облегченный вариант GAN+GA без тяжелых фреймворков машинного обучения: нейросетевые компоненты написаны на NumPy и используются для разнообразия популяции и модификации fitness.

Преимущество генетического алгоритма заключается в способности исследовать множество разных порядков размещения. Для 3D Bin Packing это существенно, потому что небольшое изменение первых элементов очереди может полностью изменить дальнейшую структуру укладки. Операторы скрещивания и мутации позволяют комбинировать удачные фрагменты разных решений, а турнирный отбор поддерживает давление в сторону более качественных особей.

Однако генетический подход чувствителен к способу кодирования. Если хромосома задает только порядок коробов, алгоритм может долго искать подходящие ориентации. Если она кодирует и порядок, и ориентации, пространство поиска становится больше, а случайные мутации чаще портят работоспособные комбинации. Поэтому в проекте выбран компромисс: генетическая часть управляет перестановкой и выбором ориентаций, а проверка геометрии остается в декодере и валидаторе.

ML-компонента в такой постановке не заменяет решатель целиком. Она используется как механизм формирования более разнообразных и потенциально перспективных кандидатов. Это важно для исследовательской работы: можно сравнить, дает ли простая обучаемая надстройка выигрыш по сравнению с классической метаэвристикой. В текущих результатах GAN+GA демонстрирует конкурентный уровень, но уступает LNS, что указывает на необходимость более богатого обучения или более тесной связи модели с геометрическими признаками задачи.

\subsection{\centering{Точный перебор и branch-and-bound}}

Для малых экземпляров возможен точный или почти точный поиск. В проекте каталог \code{brute\_force\_vs\_algo} содержит решатель branch-and-bound с backtracking и ограниченным перебором по extreme points. Для задач до 12 объектов используется полный режим, для более крупных -- limited discrepancy search с ограниченным числом отклонений от жадного порядка.

Такая реализация важна не как основной промышленный решатель, а как исследовательская точка сравнения. Она позволяет на малых наборах оценить, насколько эвристические решения близки к результатам более полного поиска. Сохраненные результаты показывают, что на малых задачах LNS близок к brute force по итоговому score, но работает существенно быстрее на более сложных наборах.

Branch-and-bound строит дерево вариантов, где очередной уровень соответствует выбору позиции и ориентации для следующего объекта. Если частичное решение уже не может привести к результату лучше текущего лучшего, соответствующая ветвь отсекается. В задачах упаковки такие оценки могут учитывать оставшийся объем, число неразмещенных объектов, допустимую высоту и грубую верхнюю границу будущего score. Даже простые отсечения сильно уменьшают перебор, но не устраняют экспоненциальный рост в худшем случае.

Limited discrepancy search использует другое наблюдение: жадный порядок часто близок к хорошему решению, но иногда требует нескольких исправлений. Вместо полного перебора всех перестановок метод ограничивает число отклонений от базового жадного выбора. Это дает промежуточный режим между greedy и полным branch-and-bound: качество повышается, но число ветвей остается управляемым для средних тестовых примеров.

В рамках данной работы точные и почти точные методы нужны прежде всего для верификации масштаба ошибки эвристик. Если LNS на малых задачах дает результат, близкий к перебору, это повышает доверие к его поведению на больших экземплярах, где точный расчет уже невозможен. Если же на малых наборах эвристика систематически проигрывает, значит проблема связана не только с размером задачи, но и с самой логикой построения решения.

\subsection{\centering{Выводы по разделу}}

Рассмотренные методы различаются по компромиссу между скоростью и качеством. Greedy и MaxRects быстры и интерпретируемы, но ограничены локальностью или послойной структурой. LNS добавляет возможность перестройки уже найденной укладки и поэтому лучше подходит для сложных экземпляров. GAN+GA расширяет поиск за счет эволюционных и генеративных механизмов, но требует настройки и не гарантирует превосходства над хорошо подобранной метаэвристикой. Brute force полезен для анализа малых задач, но не масштабируется на рабочие размеры.

Для проекта Smart 3D Packing наиболее важен не отдельный алгоритм сам по себе, а сопоставимость методов в единой постановке. Все подходы должны работать с одинаковым входным JSON, возвращать размещения в едином формате и проверяться одним валидатором. Только в этом случае сравнение score, утилизации объема, полноты заказа, хрупкости и времени выполнения отражает свойства алгоритмов, а не различия в правилах проверки.
